\documentclass{article}
\usepackage[T1]{fontenc}
\usepackage{lmodern}
\usepackage{graphicx}
\usepackage{amsmath,amssymb}
\usepackage[a4paper,margin=2cm]{geometry}
\usepackage[absolute,overlay]{textpos}
\setlength{\TPHorizModule}{1mm}
\setlength{\TPVertModule}{1mm}
\usepackage[indonesian]{babel}
\usepackage{hyperref}
\setlength{\emergencystretch}{2em}
% unit: Halaman 26

\begin{document}

% === Halaman 26 (llm) ===

\{ Dekripsi RC6-w/r/b \\
Masukan: Cipherteks disimpan di dalam empat register berukuran $w$ bit: $A, B, C$ dan $D$; $r$ adalah jumlah putaran; $S[0, \dots, 2r+3]$ adalah kunci internal, panjangnya $w$ bit \\
Luaran: Plainteks disimpan di dalam $A, B, C, D$
\}

Algoritma:\
\[ C \leftarrow C - S[2r+3] \\
A \leftarrow A - S[2r+2] \]

\begin{algorithmic}
\For{$i \leftarrow r \text{ downto } 1$}
\State $(A, B, C, D) \leftarrow (D, A, B, C)$
\State $u \leftarrow (D^*(2D + 1)) \lllg lg w$
\State $t \leftarrow (B^*(2B + 1)) \lllg lg w$
\State $C \leftarrow ((C - S[2i + 1]) \gggt t) \oplus u$
\State $\mathbf{A \leftarrow ((A - S[2i]) \gggt u) \oplus t}$
\EndFor
\end{algorithmic}
\[ D \leftarrow D - S[1] \\
B \leftarrow B - S[0] \]

\end{document}
